\thispagestyle{empty}

\subsection*{Abstract}

The Dataspace ecosystem is growing as a reliable solution for organizations to securely and sovereignly share data within a decentralized framework, with EDCs being the backbone of this system. However, performance analysis and benchmarking for these systems is still a gap. This thesis aims to conduct a systematic performance benchmark, mainly focusing on four variants of EDCs. Furthermore, the thesis will also analyze the results to identify bottlenecks and areas for improvement in EDC implementations, providing insights for future research and development in this field.

The research will adapt the GQM (Goal-Question-Metric) approach to define the performance goals, formulate research questions, and identify relevant metrics for evaluation. The study will involve setting up a controlled environment to conduct experiments and collect data on various performance aspects of EDCs, such as latency, throughput, resource utilization, and scalability under different workloads and configurations.

The results of this thesis will contribute to the understanding of the performance characteristics of EDCs and provide valuable insights for practitioners and researchers in the field of data exchange and integration.

\noindent  \\
